\section{Alcance, organización y trazabilidad}
El plan cubre E1, E2, sus marchas blancas y las liberaciones de los 36 meses de Operación. Calidad coordinará el expediente; Desarrollo ejecutará y corregirá pruebas de construcción; Datos comprobará migración y conciliación; Seguridad verificará controles e intrusión independiente; Operación comprobará despliegue, continuidad y retorno. La Contraparte Técnica validará escenarios de negocio y decidirá aceptación conforme a T-17. Se conservan responsables y capacidad de T-14/T-15, sin agregar una estimación duplicada. \parencites[sec.~9.1--9.3]{sd9-sd9}[Diccionario de gobierno, pruebas, formación y servicio]{sd9-t14}

Cada caso conservará ID propio, ID y cláusula de T-12, etapa, hito, paquete EDT, diseño/contrato, condición inicial, datos, acciones, resultado esperado, umbral, ambiente, automatización y responsable. Al ejecutarse añadirá fecha efectiva, commit/digest, configuración, resultado observado, evidencia y defectos. Los casos se versionarán con el requisito y el artefacto; una modificación invalidará los resultados afectados. La matriz conservará los 576 IDs/enunciados y los 27 complementos específicos. Antes de aceptar un alcance se exigirá cobertura del 100\,\% de requisitos aplicables con todas sus cláusulas verificadas; una prueba sin ejecutar sólo acredita planificación.

El expediente aplicará ISO/IEC/IEEE 29119-2:2021 a la organización y ejecución de pruebas, 29119-3:2021 a sus registros y 29119-4:2021 al diseño de casos. Calidad conservará la correspondencia entre procesos, documentos y técnicas seleccionadas, con responsable y evidencia por versión. Incluirá plan, diseño y datos, preparación del ambiente, registro de ejecución, incidencias, estado de avance e informe de cierre; una ficha de caso aislada no acreditará la conformidad de todo el proceso. \parencite{sd9-29119-2,sd9-29119-3,sd9-29119-4}

\section{Niveles, entradas y salidas}
La Tabla~\ref{tab:t13-niveles} fija las condiciones de entrada y salida por nivel. Las salidas incorporan los umbrales completos de SD9 9.1, no sólo el indicador resumido aquí. Un control ausente, fallido o correspondiente a otra versión bloqueará la promoción. \parencites[20.1, pp.~34--35]{sd9-bt}[sec.~9.1--9.2]{sd9-sd9}
\begin{longtable}{L{.19\textwidth}L{.36\textwidth}L{.36\textwidth}}
\caption{Condiciones por nivel de prueba}\label{tab:t13-niveles}\\\hline
\EncabezadoTabla{Nivel} & \EncabezadoTabla{Entrada} & \EncabezadoTabla{Salida}\\\hline\endfirsthead
\hline\EncabezadoTabla{Nivel} & \EncabezadoTabla{Entrada} & \EncabezadoTabla{Salida}\\\hline\endhead
Unidad/componente & Regla, cambio y datos identificados. & Cero fallas; cobertura de negocio $\geq70\,\%$ por componente.\\\hline
Integración & Artefacto recibido en QA y contratos versionados. & Todos los flujos exigibles sin error ni efecto duplicado.\\\hline
Sistema/regresión & Versión integral, perfiles y dependencias disponibles. & Batería obligatoria completa sin regresiones.\\\hline
Aceptación de usuario & Ejecuciones técnicas completas y conformes; usuarios preparados. & Casos obligatorios aprobados y firmados por la Contraparte.\\\hline
\end{longtable}
\FuenteTabla{Bases Técnicas Transversales, 20.1; criterios de Synaptix en SD9.}

Las pruebas aplicarán particiones, límites, decisiones y transiciones de estado según la regla evaluada. Procesos, registros y técnicas siguen ISO/IEC/IEEE 29119-2/3/4:2021. El resultado esperado procederá del requisito y contrato, sin copiar el resultado del programa como patrón de corrección. \parencite{sd9-29119-2,sd9-29119-3,sd9-29119-4}

\section{Tipos y escenarios obligatorios}
La Tabla~\ref{tab:t13-tipos} identifica controles especiales. Cada tipo se ejecutará sobre el alcance completo aplicable a la versión y dejará protocolo, registros brutos, informe y decisión. Las pruebas que cambien carga, topología o datos se separarán en ambientes aislados o se secuenciarán con restauración verificada. \parencite[20.1, pp.~34--35]{sd9-bt}
\begin{longtable}{L{.20\textwidth}L{.36\textwidth}L{.35\textwidth}}
\caption{Escenarios especiales y criterio de salida}\label{tab:t13-tipos}\\\hline
\EncabezadoTabla{Tipo} & \EncabezadoTabla{Escenario} & \EncabezadoTabla{Salida}\\\hline\endfirsthead
\hline\EncabezadoTabla{Tipo} & \EncabezadoTabla{Escenario} & \EncabezadoTabla{Salida}\\\hline\endhead
Carga/rendimiento & Mezcla conjunta a $1,5$ veces peak, con historia y retención representativas. & Todos los umbrales P95 y máximos propios de SD9 cumplidos.\\\hline
Estrés & Aumento progresivo hasta saturación y posterior reducción. & Curva respuesta/carga, recursos, saturación/quiebre, degradación controlada y recuperación sin pérdida silenciosa.\\\hline
Resiliencia & Caída de instancia/zona, enlace o dependencia externa; latencia elevada y saturación de disco. & Degradación controlada y recuperación sin intervención, según escenario.\\\hline
Desastre & Conmutación regional real, identidad, datos y retorno DR-4. & RTO $\leq4$ h y RPO $\leq15$ min, con conciliación.\\\hline
Desconexión & Recinto 24 h; terreno 12 h; boyas 8 h; reconexión con registros nuevos. & Funciones exigibles completas; reconexión $\leq30$ min sin pérdida ni duplicados.\\\hline
Seguridad ofensiva & Intrusión independiente sobre superficie y perfiles completos. & Cero hallazgos críticos/altos abiertos y correcciones verificadas.\\\hline
Accesibilidad & Comprobación automática y manual por tarea, teclado y dispositivo. & Conformidad WCAG 2.2 AA documentada.\\\hline
Migración & Dos ensayos completos válidos de extracción a retorno por etapa. & Conciliación sin diferencias no explicadas y controles críticos de SD5 conformes.\\\hline
Despliegue/retorno & Versiones N/N+1, conexiones en vuelo y hechos nuevos. & Sin interrupción ni pérdida, reversa automatizada y autoridad íntegra.\\\hline
Usabilidad & Tareas por perfil con usuarios y dispositivos representativos. & Tiempos, interacciones, error y aprendizaje de SD9 9.1.4 cumplidos.\\\hline
\end{longtable}
\FuenteTabla{Bases Técnicas Transversales, cap.~20; C07, cap.~15; SD4 DR-4; SD5 y SD9.}

Los escenarios funcionales recorrerán navegación y zarpes, amarras, escuela, combustible y consumos, inspección de boyas, varadero, contratistas, administración y hechos facturables, y gestión ambiental, con los portales, avisos y servicios comunes de SD3--SD5. Se probarán permisos cruzados denegados, autorizaciones vencidas/revocadas, pagos de resultado incierto, despachos repetidos, mensajes fuera de orden y continuidad de la autoridad local. La situación escolar compartida y el saldo no se validarán con una copia periódica como sustituto de tiempo real. Las innovaciones conservarán además sus pruebas específicas y la integración del conjunto. \parencites[cap.~15, pp.~34--37]{sd9-c07}[sec.~5.2--5.4]{sd9-sd5}

Los controles de seguridad operacional de SD9 se ejecutarán en CP-04, CP-06 y CP-07 con fallas, restricciones y alertas identificadas. CP-08 incluirá reglas de captura y saneamiento, indicadores de completitud y exactitud, contraste con fuentes autorizadas y registro de correcciones; CP-17 comprobará actualidad. Datos conservará los denominadores y la cobertura de contraste del universo exigible. La recepción aplicará los criterios completos de SD9, con evidencia independiente de las pruebas de código. \parencites{sd9-25012}[sec.~9.1.4]{sd9-sd9}

\section{Ambientes, datos y automatización}
Desarrollo y QA comprobarán código e integración; Preproducción sostendrá las pruebas completas con EQ-4; Producción y Recuperación ante Desastres conservarán aislamiento y autoridades para las comprobaciones reales autorizadas. La preparación verificará topología, versiones, configuración, capacidad, cuotas y disponibilidad de los equipos. Se ensayará crecimiento con datos y almacenamiento representativos, no sólo con sesiones adicionales. \parencite[sec.~6.2]{sd9-sd6}

Los datos no productivos serán sintéticos o anonimizados/seudonimizados verificablemente, con origen autorizado, manifiesto, integridad, versión y regla de conservación/eliminación. Los secretos serán propios del ambiente. Los avisos, cargos y otros efectos externos usarán destinos de ensayo autorizados; una simulación no acreditará la disponibilidad efectiva de la integración. Se conservarán registros fallidos, sin descartar abandonos o tiempos agotados del análisis. \parencites[RT-11.25, p.~24]{sd9-bt}[sec.~9.2.3]{sd9-sd9}

GitHub Actions ejecutará pruebas y controles versionados con SonarQube, Trivy, Gitleaks y ZAP conforme a SD6, conservando reportes, SBOM, firma y procedencia verificadas con el digest promovido. La promoción será del mismo artefacto sin recompilar y el retorno estará automatizado. Las pruebas manuales y la intrusión independiente conservarán resultados propios; una etapa omitida no se tratará como aprobada. \parencite[sec.~6.2]{sd9-sd6}

La Tabla~\ref{tab:t13-ejecucion} integra ambientes, datos y automatización por nivel y tipo; sus entradas y salidas se completan con las Tablas~\ref{tab:t13-niveles} y \ref{tab:t13-tipos}. En todos los casos la entrada exige requisito, protocolo, versión, datos y ambiente disponibles; la salida exige resultado conforme, evidencia íntegra y bloqueos cerrados.
\begin{longtable}{L{.22\textwidth}L{.22\textwidth}L{.22\textwidth}L{.24\textwidth}}
\caption{Matriz de ejecución, datos y automatización}\label{tab:t13-ejecucion}\\\hline
\EncabezadoTabla{Nivel/tipo} & \EncabezadoTabla{Ambiente} & \EncabezadoTabla{Datos} & \EncabezadoTabla{Automatización}\\\hline\endfirsthead
\hline\EncabezadoTabla{Nivel/tipo} & \EncabezadoTabla{Ambiente} & \EncabezadoTabla{Datos} & \EncabezadoTabla{Automatización}\\\hline\endhead
Unidad/componente & DEV y CI & Sintéticos y límites & GitHub Actions, cada cambio; cobertura y cero fallas.\\\hline
Integración & QA, cada despliegue & Sintéticos; terceros de ensayo & Contratos OpenAPI/AsyncAPI y todos los flujos automatizados.\\\hline
Sistema/regresión & QA, antes de PREPROD & Conjuntos sintéticos o anonimizados & Batería completa en GitHub Actions.\\\hline
UAT y usabilidad & PREPROD & Casos por perfil protegidos & Ejecución asistida y medición; firma real del CLIENTE.\\\hline
Carga/estrés & PREPROD equivalente & Mezcla y volumen de SD5 & Generador y métricas versionados; informe de curva, recursos y saturación.\\\hline
Resiliencia & PREPROD; PROD autorizado en Operación & Hechos controlados y dependencias & Inyección y comprobaciones repetibles; supervisión de Operación.\\\hline
DRP & Primario/DR reales, con autorización & Hechos identificados, réplicas y conciliación & Conmutación, medición y retorno; validación funcional supervisada.\\\hline
Seguridad & CI, QA y PREPROD & Identidades por rol y datos protegidos & SonarQube/Trivy/Gitleaks/ZAP; intrusión por tercero independiente.\\\hline
Accesibilidad & QA y PREPROD & Tareas de todas las interfaces & Comprobación automática y manual WCAG 2.2 AA.\\\hline
Migración & PREPROD & Universo exigible protegido de SD5 & Dos ensayos completos; extracción a retorno, conciliación automatizada y revisión de Datos.\\\hline
\end{longtable}
\FuenteTabla{Synaptix, síntesis de SD6, SD9 y Bases Técnicas Transversales, caps.~4 y 20.}

Cada contrato OpenAPI/AsyncAPI se vinculará con sus operaciones, esquemas, mensajes y casos positivos/negativos; QA ejecutará el 100\,\% de flujos aplicables sin error tras cada despliegue. La matriz de trazabilidad registrará ID y cláusula de T-12, etapa/hito, decisión y ubicación de diseño en SD4/SD5, módulo/ruta/commit de código, ID de caso, conjunto de datos, ejecución/resultado/evidencia, defecto y versión/digest/ambiente/fecha de despliegue LIB-6. Permitirá buscar en ambos sentidos y distinguir caso diseñado de caso ejecutado conforme. Se conservan los 576 IDs/enunciados y los 27 complementos. \parencite[sec.~9.2.4]{sd9-sd9}

La seguridad conservará el mapa de controles OWASP ASVS 4.0 nivel 2 como mínimo, con identificadores de la revisión 4.0.3, evidencias y exclusiones justificadas. La intrusión independiente se ejecutará antes de cada paso a producción y anualmente en Operación. No se permitirá promoción con críticos/altos abiertos; remediación máxima de siete/quince/treinta días corridos para críticas/altas/medias desde publicación o detección. \parencites{sd9-asvs}[RT-11.04/20; sec.~11.5]{sd9-bt}

\section{Calendario y dependencias de aceptación}
UAT E1 ocupará del 12 al 15 de octubre de 2027, revisión del 18 al 29 y subsanación del 2 al 15 de noviembre. H5 se programará el 16 de noviembre, después del cierre de observaciones de los ensayos; la reserva marina será del 16 al 29 antes de activar el 1 de diciembre. Se conservan 48 HH y los cuatro recursos de la asignación integral T-15. La programación obliga a repetir evidencia invalidada por cambios y no anticipa actas de conformidad. \parencites{sd9-t15}[arts.~17--18]{sd9-ba}

El calendario adopta inicio el 1 de diciembre de 2026 y las asignaciones de T-15. La Tabla~\ref{tab:t13-calendario-inicial} conserva las ventanas iniciales de construcción y certificación de cada etapa. Las pruebas iniciales no sustituyen las verificaciones antes de cada paso a producción, incluidas las nuevas liberaciones durante Operación; una versión o perfil cambiado repetirá las pruebas afectadas. \parencites[sec.~7.3]{sd9-sd7}{sd9-t15}[20.1, p.~34]{sd9-bt}
\begin{longtable}{L{.28\textwidth}L{.31\textwidth}L{.31\textwidth}}
\caption{Ventanas iniciales de pruebas integrales}\label{tab:t13-calendario-inicial}\\\hline
\EncabezadoTabla{Tipo / paquete} & \EncabezadoTabla{E1} & \EncabezadoTabla{E2}\\\hline\endfirsthead
\hline\EncabezadoTabla{Tipo / paquete} & \EncabezadoTabla{E1} & \EncabezadoTabla{E2}\\\hline\endhead
Carga/estrés, \texttt{1.10.2.e.4} & 24--26 agosto 2027 & 16--20 marzo 2028\\\hline
Resiliencia, \texttt{1.10.2.e.6} & 1--8 septiembre 2027 & 24--31 marzo 2028\\\hline
DR, \texttt{1.10.2.e.7} & 8--13 septiembre 2027 & 31 marzo--5 abril 2028\\\hline
Ofensiva, \texttt{1.10.2.e.8} & 24--30 agosto 2027 & 16--22 marzo 2028\\\hline
\end{longtable}
\FuenteTabla{Synaptix, programación integrada de T-15 y SD7.}

Carga y seguridad con fechas coincidentes requerirán instancias aisladas y recursos autorizados; en caso contrario deberán secuenciarse. Las fechas no dispensan la restauración ni la independencia del ensayo. UAT exigirá ejecuciones técnicas completas y resultados conformes de la misma versión; la revisión contractual de sus expedientes podrá continuar en paralelo. La certificación esperará el cierre de todas las observaciones impeditivas, conservando revisión y subsanación del artículo 18. Un cambio derivado de esa revisión obligará a repetir la verificación afectada, incluida UAT cuando corresponda. H5/H10 y H7/H12 mantendrán sus meses contractuales, por lo que Proyecto comprobará capacidad y cabida antes de liberar la agenda final. \parencite[art.~18, p.~13]{sd9-ba}

Los ensayos completos de migración conservarán las ventanas y puertas VM/VS de SD9 9.3 y T-14; cada ensayo reserva diez días hábiles de ejecución, diez de revisión y diez de subsanación. La marcha blanca conserva observación real e indicadores diarios de T-18. Durante Operación se aplicará el calendario periódico de la Tabla~\ref{tab:t13-periodico}; los controles antes de cada producción se añadirán al cambio correspondiente, sin quedar absorbidos por una fecha periódica. \parencites[Ensayos completos y recepción previa a los cortes]{sd9-t14}[sec.~9.3]{sd9-sd9}
\begin{longtable}{L{.22\textwidth}L{.27\textwidth}L{.41\textwidth}}
\caption{Calendario periódico de pruebas en Operación}\label{tab:t13-periodico}\\\hline
\EncabezadoTabla{Control} & \EncabezadoTabla{Paquete} & \EncabezadoTabla{Intervalo programado}\\\hline\endfirsthead
\hline\EncabezadoTabla{Control} & \EncabezadoTabla{Paquete} & \EncabezadoTabla{Intervalo programado}\\\hline\endhead
Resiliencia & \texttt{1.12.4.1.6} & 2028-11-03 a 2028-11-09\\\hline
DR y retorno & \texttt{1.12.4.1.7} & 2028-11-07 a 2028-11-10\\\hline
Resiliencia & \texttt{1.12.4.2.6} & 2029-03-21 a 2029-03-27\\\hline
DR y retorno & \texttt{1.12.4.2.7} & 2029-03-21 a 2029-03-27\\\hline
Ofensiva anual & \texttt{1.12.5.1} & 2029-03-19 a 2029-03-26\\\hline
Resiliencia & \texttt{1.12.4.3.6} & 2029-08-02 a 2029-08-08\\\hline
DR y retorno & \texttt{1.12.4.3.7} & 2029-08-06 a 2029-08-09\\\hline
Resiliencia & \texttt{1.12.4.4.6} & 2029-11-06 a 2029-11-12\\\hline
DR y retorno & \texttt{1.12.4.4.7} & 2029-11-08 a 2029-11-13\\\hline
Resiliencia & \texttt{1.12.4.5.6} & 2030-03-21 a 2030-03-28\\\hline
DR y retorno & \texttt{1.12.4.5.7} & 2030-03-21 a 2030-03-27\\\hline
Ofensiva anual & \texttt{1.12.5.2} & 2030-03-19 a 2030-03-26\\\hline
Resiliencia & \texttt{1.12.4.6.6} & 2030-08-02 a 2030-08-08\\\hline
DR y retorno & \texttt{1.12.4.6.7} & 2030-08-06 a 2030-08-09\\\hline
Resiliencia & \texttt{1.12.4.7.6} & 2030-11-05 a 2030-11-11\\\hline
DR y retorno & \texttt{1.12.4.7.7} & 2030-11-07 a 2030-11-12\\\hline
Resiliencia & \texttt{1.12.4.8.6} & 2031-03-20 a 2031-03-27\\\hline
DR y retorno & \texttt{1.12.4.8.7} & 2031-03-20 a 2031-03-26\\\hline
Ofensiva anual & \texttt{1.12.5.3} & 2031-03-18 a 2031-03-25\\\hline
Resiliencia & \texttt{1.12.4.9.6} & 2031-07-02 a 2031-07-08\\\hline
DR y retorno & \texttt{1.12.4.9.7} & 2031-07-04 a 2031-07-09\\\hline
\end{longtable}
\FuenteTabla{Synaptix, programación integrada de T-15; frecuencias mínimas de BT, cap.~20.}


\section{Fichas de escenarios críticos y datos}
Las fichas siguientes fijan condiciones iniciales, acciones y resultados verificables para los principales riesgos de la marina. Sus IDs identifican casos de T-13 y no reemplazan los IDs ni enunciados de T-12. La aplicación a E1/E2 dependerá del alcance habilitado de cada requisito; antes de ejecutar se cerrará el manifiesto de datos, versión y cláusulas cubiertas. Los resultados y actas se producirán al realizar cada prueba. \parencites[cap.~15, pp.~34--37]{sd9-c07}[cap.~5, RT-05.13/14, pp.~11--13; cap.~20, pp.~34--35]{sd9-bt}[sec.~9.1--9.2]{sd9-sd9}
\subsection{CP-01: operación del recinto desconectado}
\label{caso:cp-01}
\textbf{Trazabilidad:} C07 RT-03.10; SD4/SD5; \texttt{1.10.2.e.5}. Datos preparará el manifiesto inicial de embarcaciones, amarras, permisos, escuela, stock y registros vigentes contra su autoridad, además de operaciones nuevas para cada función local exigible. Se cortarán todos los enlaces exteriores durante 24 horas completas, manteniendo la carga y operación autorizadas del recinto; Calidad comprobará recepción y asignación de embarcaciones, salida/regreso, zarpe, combustible y escuela. Se exigirá ejecución electrónica durante todo el horizonte, persistencia durable, cero efectos duplicados y consulta coherente de autoridad. El diario manual se registrará como contingencia y no como función electrónica cumplida. Evidencia: maniobra de corte, energía/conectividad, registros temporales y conciliación por función.

\subsection{CP-02: terreno y boyas fuera de cobertura}
\label{caso:cp-02}
\textbf{Trazabilidad:} C07 RT-03.10; \texttt{1.10.2.e.5}. Se prepararán dispositivos autorizados de pantalán, varadero, escuela y embarcación de inspección con conjuntos propios del perfil y registro de batería. Se ejecutarán tareas y hechos nuevos durante 12 horas continuas de terreno y ocho horas de inspección, incluyendo reinicio del dispositivo y almacenamiento de evidencia autorizado. Se exigirá conservación de cada registro y su identidad, cero duplicados al reconectar y protección de datos locales. La capacidad se comprobará sobre la jornada completa; un dispositivo con cola vacía no acredita el horizonte. Se registrará separadamente cualquier función que requiera autoridad compartida para no equiparar captura local con decisión global disponible.

\subsection{CP-03: sincronización después de 24 horas}
\label{caso:cp-03}
\textbf{Trazabilidad:} C07 RT-03.13; SD5 5.4; \texttt{1.10.2.e.5}. Se utilizará el universo acumulado en CP-01, con nuevas operaciones concurrentes, duplicados de transporte y mensajes fuera de orden. Tras restituir el enlace, se medirá desde reconexión hasta recepción, aplicación durable, actualización y conciliación completa. Se exigirá duración máxima de treinta minutos, cero pérdida de salidas, regresos, despachos o asistencia y cero efectos repetidos. El manifiesto incluirá tamaño, cantidad, caudal efectivo y hechos nuevos; una estimación de transferencia no sustituirá el tiempo total observado. Datos y Calidad comprobarán cada ID y diferencias; no se purgará la cola para obtener el umbral.

\subsection{CP-04: situación escolar y permisos revocados}
\label{caso:cp-04}
\textbf{Trazabilidad:} C07 RT-05.29; SD5, autoridad escolar. Se prepararán alumnos, apoderados, monitores, botes y permisos vigentes, vencidos y revocados, con datos protegidos. Se registrarán salida, cambio, regreso y retiro desde dos puestos autorizados, incluyendo partición de comunicación y reintento. Se exigirá reflejo de cada hecho contra la autoridad compartida en tiempo real, rechazo de actuación sin autorización y cero acceso cruzado entre apoderados. Una partición que impida conocer quién está en el agua se registrará como incumplimiento aunque exista copia local. Evidencia: secuencia de hechos, consulta de ambos puestos y decisiones de permisos.

\subsection{CP-05: tiempos completos de operación}
\label{caso:cp-05}
\textbf{Trazabilidad:} C07 RT-09.01; \texttt{1.10.2.e.4/10}. Funcional preparará embarcación visitante compatible y no compatible, registros de salida/regreso, personas y antecedentes de zarpe y clase de hasta 18 botes. Usuarios representativos ejecutarán las tareas con dispositivos y condiciones de uso seguras del perfil. Se medirán consulta radial a confirmación de amarra $\leq45$ s; inicio a persistencia de salida/regreso $\leq10$ s y máximo dos interacciones; nómina/legajo $\leq90$ s; y clase completa $\leq60$ s. Se conservarán tiempos individuales y errores; medir sólo API no acredita la tarea. La aceptación exigirá resultados correctos dentro de cada máximo.

\subsection{CP-06: avisos de emergencia y planes vencidos}
\label{caso:cp-06}
\textbf{Trazabilidad:} C07 RT-09.01 y RT-16.21; módulo ALE. Se prepararán 296 destinatarios de ensayo con mensaje de texto, mensajería instantánea y correo, contactos válidos, errores de entrega y acuses ausentes, más planes sin cierre con vencimiento identificado. Se activará una emergencia y se comprobarán los tres canales simultáneos, reintentos y escalamiento a llamada registrada para quienes no acusen; los destinos de ensayo evitarán avisos reales no autorizados. El último envío deberá ocurrir dentro de diez minutos, registrando aparte entregas y acuses por destinatario. La alerta de plan vencido se generará en máximo quince minutos desde vencimiento, sin agregar otro período de gracia. Evidencia: destinatarios, intentos, respuestas, pendientes y escalamiento.

\subsection{CP-07: combustible y efectos externos idempotentes}
\label{caso:cp-07}
\textbf{Trazabilidad:} SD5, autoridad local y transacciones; C07 RT-03.10/13. Se prepararán operaciones de despacho y hechos facturables con ID y versión, confirmaciones repetidas y respuestas externas inciertas. Se repetirán solicitudes, se interrumpirá comunicación después de persistir y antes del acuse y se reiniciará el productor. Se exigirá un único efecto por hecho y correspondencia entre movimiento, evidencia y resultado externo. Una respuesta incierta se consultará y conciliará según el contrato; no se reintentará como fallo definitivo si eso pudiera duplicar un cargo. Evidencia: IDs, transacciones, inbox/outbox, acuses y conciliación.

\subsection{CP-08: migración y retorno con hechos nuevos}
\label{caso:cp-08}
\textbf{Trazabilidad:} BT RT-05.13/14, cap.~20; \texttt{1.10.1.e.j.k}. Datos preparará el universo histórico exigible de SD5, incluidos los catorce expedientes de abandono, relaciones y documentos asociados, con manifiesto protegido y reglas versionadas. Cada ensayo recorrerá extracción, transformación, carga, conciliación y retorno, incluyendo hechos nuevos posteriores al corte y resultados externos inciertos. Se exigirá cero diferencias no explicadas y las tolerancias críticas de SD5, sin compensar saldos de signo contrario ni descartar hechos posteriores al restaurar. Ambos ensayos deberán ser completos y válidos para la versión definitiva; el segundo incluirá suplentes técnico y funcional. Evidencia: manifiestos, tiempos por fase, diferencias y decisión de retorno.

\subsection{CP-09: conmutación regional y retorno DR-4}
\label{caso:cp-09}
\textbf{Trazabilidad:} BT RT-07.01--06; \texttt{1.10.2.e.7} y \texttt{1.12.4}. Operación preparará versión, réplica, identidades, claves, objetos, colas y corte de hechos confirmados en los sitios comprometidos. Se aislará el origen, se promoverá el destino y se comprobará capacidad, autenticación, datos y operaciones antes del retorno. Se exigirá RTO $\leq4$ h desde interrupción hasta funcionalidad restituida y RPO $\leq15$ min sobre hechos regionales confirmados, sin doble escritor ni efectos repetidos. Se distinguirán hechos de terreno aún no recibidos y su obligación local. Evidencia: cronología, último dato confirmado, permisos, fencing, conciliación y retorno; una réplica existente no acredita el ensayo.

\subsection{CP-10: acceso cruzado y revocación}
\label{caso:cp-10}
\textbf{Trazabilidad:} BT cap.~11; SD4, identidad y segregación; \texttt{1.10.2.e.8}. Seguridad preparará cuentas de los perfiles comprometidos, sesiones emitidas, credenciales revocadas y solicitudes a recursos propios y ajenos. Se ejecutarán accesos negativos por interfaz directa y portal, revocación durante sesión y comprobación de registros sin exposición innecesaria de información personal. Se exigirá cero lectura/escritura no autorizada y cierre de todos los hallazgos críticos/altos. La intrusión independiente conservará alcance y reporte íntegro; el análisis automático complementa esa comprobación y no la reemplaza.

\subsection{CP-11: accesibilidad y aprendizaje por perfil}
\label{caso:cp-11}
\textbf{Trazabilidad:} BT RT-13.01--11; \texttt{1.10.2.e.9/10}. Se prepararán tareas por perfil, dispositivos soportados, navegación de teclado y ayudas técnicas, además de sesiones de instrucción registradas. Se comprobarán criterios WCAG 2.2 AA aplicables automática y manualmente, foco, errores, objetivos táctiles y uso seguro con guantes. Se medirán tiempo, interacciones, ayuda, errores y aprendizaje con las fórmulas de SD9, sin promediar perfiles ni borrar recuperaciones. Se exigirá 95\,\% de tareas sin asistencia, error $\leq5\,\%$, actuaciones críticas 100\,\% correctas y objetivo de interfaz tras cuatro horas de instrucción; la formación adicional se registrará como recuperación.

\subsection{CP-12: promoción y reversa sin interrupción}
\label{caso:cp-12}
\textbf{Trazabilidad:} BT RT-04.06/07/10; DEP-6/MIG-6; SD6. Se prepararán artefactos anterior/nuevo, digest, configuración, esquema compatible y solicitudes/conexiones en vuelo sobre EQ-4 a $1,5$ veces peak. Se ejecutarán promoción y retorno automatizados, con falla inducida y hechos nuevos durante el cambio. Se exigirá cero interrupción, pérdida o duplicación, cumplimiento de tiempos y conservación de autoridad y sesión. El ensayo comprobará expandir/transitar/retirar esquema y retorno representable. Una conmutación de infraestructura cuya continuidad no esté demostrada permanecerá bloqueada; una ventana o RTO conforme no dispensará ese criterio.

\subsection{CP-13: crecimiento por parametrización}
\label{caso:cp-13}
\textbf{Trazabilidad:} C07 RT-02.12; SD4 y modelo de SD5. Se preparará una copia de ensayo del catálogo con pantalanes, amarras, puestos de invernada y boyas, manteniendo sus relaciones y reglas de compatibilidad. Una persona usuaria autorizada del perfil de administración configurará nuevos elementos hasta el escenario de 380 amarras y 70 boyas, usando funciones ordinarias y documentación del producto. Se exigirá cero cambio de código o rediseño, cero intervención de especialista para parametrizar y funcionamiento correcto de asignación, búsqueda, permisos, telemetría y consultas sobre las unidades nuevas. El incremento de recursos de capacidad se distinguirá del acto de parametrización y deberá conservar el plan de SD4. Evidencia: parámetros y auditoría del cambio, relaciones, resultados y recursos utilizados.

\subsection{CP-14: cobertura, segregación y enlace de respaldo}
\label{caso:cp-14}
\textbf{Trazabilidad:} C07 RT-03.24. Se preparará el protocolo de recorrido y puntos de medición de los cinco pantalanes y el varadero, con el rango de marea real, flexión de cableado y condiciones seguras de desplazamiento. Se comprobará captura y consulta en cada zona, segregación efectiva entre socios, administración, cámaras y operación, y conmutación al enlace de respaldo contratado. Se exigirá cobertura funcional de todas las zonas, rechazo de accesos no autorizados entre segmentos y continuidad conforme al requisito del servicio durante la falla del enlace principal. El expediente incluirá mediciones, diagrama real, pruebas negativas y evidencia del contrato y recepción del respaldo, sin atribuirlo al equipamiento existente. Una visita en una única condición de marea no cerrará la comprobación del rango.

\subsection{CP-15: retención, custodia y eliminación}
\label{caso:cp-15}
\textbf{Trazabilidad:} C07 RT-05.10; SD5. Se prepararán registros de cada categoría con fechas de origen y vencimiento, originales y copias, bloqueos legales de eliminación y permisos de acceso. El ensayo verificará con tiempo controlado los límites: menores hasta 18 años más cinco y mínimo diez años; salida/regreso y zarpe cinco años; contratos/custodia/cuenta seis; residuos/derrame seis; consumos cinco; fondeo vida útil más cinco; combustible cinco; video seis meses. Para menores se aplicará el vencimiento más tardío entre ambas condiciones. Antes de vencer se exigirá recuperación íntegra por perfil autorizado; después se comprobará el proceso de eliminación segura de las copias bajo su regla y el registro de la actuación, sin borrar evidencia sujeta a custodia legal. Una simulación de fecha verificará la regla, sin declararse una retención histórica ya ejecutada.

\subsection{CP-16: intercambio sectorial y disponibilidad efectiva}
\label{caso:cp-16}
\textbf{Trazabilidad:} C07 RT-05.23. Integración preparará por autoridad o servicio el formato, versión, canal, credenciales, disponibilidad confirmada y contrato de intercambio para zarpe, avisos meteorológicos/marejada, residuos peligrosos, indicadores ambientales y mensajería de baja capacidad. Se comprobarán mensajes válidos, rechazo de formato incorrecto, reintentos, latencia de horas y entrega desordenada cuando correspondan. Se exigirá correspondencia con el formato y canal realmente admitidos y persistencia/conciliación del resultado. Un simulador validará el desarrollo, pero no cerrará la disponibilidad efectiva del tercero. Si una autoridad no dispone de API, se comprobará el procedimiento seleccionado y autorizado de intercambio, sin inventar interfaz o aceptación institucional.

\subsection{CP-17: actualización de datos y mezcla de peak}
\label{caso:cp-17}
\textbf{Trazabilidad:} C07 RT-05.29 y RT-09.02; SD5 5.4. Se preparará la población y mezcla derivada de la volumetría 14.1/14.2, con mañana de fin de semana largo, consultas, avisos, telemetría, historia retenida y hechos nuevos simultáneos. El ensayo a $1,5$ veces peak conservará los denominadores y tasas de SD4--SD5. Se verificará ocupación en máximo 60 s desde evento, acumulado por amarra al menos horario, saldo y situación escolar en tiempo real contra su autoridad y consolidación ambiental diaria. El informe conservará P95 de respuesta y máximos específicos del caso, cargas, recursos, errores y saturación. Una consulta sobre proyección antigua deberá indicar su condición sin presentarse como dato actual; un promedio anual no sustituirá la concentración del peak.

\subsection{CP-18: autonomía administrada y periféricos}
\label{caso:cp-18}
\textbf{Trazabilidad:} C07 RT-06.01 y RT-17.06; SD4/T-11. Se preparará el inventario seleccionado, variantes, emplazamiento y certificados aplicables de gabinetes, medidores, sensores, lectores, despacho de combustible, estación meteorológica, control de acceso y las 38 cámaras existentes. Se comprobarán compatibilidad e integración efectiva, unidades y calibración cuando corresponda, recepción de avisos y comportamiento ante fallas. Los gabinetes nuevos de administración, pantalán y varadero acreditarán la protección exigible; la instrumentación del surtidor conservará la aptitud para su clasificación. Durante operación normal se comprobará autoadministración y ausencia de intervención humana local, distinguiendo mantenimiento o incidentes extraordinarios. Un catálogo no acredita instalación ni recepción y una prueba de escritorio no reemplazará la verificación en torreta flotante.

\subsection{CP-19: personas, perfiles y bilingüismo}
\label{caso:cp-19}
\textbf{Trazabilidad:} C07 RT-12.11/12, RT-13.08/12, RT-17.01 y RT-22.04. Se preparará la matriz de perfiles internos/externos, incluyendo armadores, socios, apoderados, visitantes, contratistas/trabajadores, charter, proveedores/gestores y autoridades aplicables. Los cinco perfiles móviles recorrerán sus tareas: pantalán, escuela, varadero, inspección y armador en baja capacidad. Se comprobarán guantes, pantalla mojada y sol directo en pruebas seguras; no se exigirá interacción durante maniobra con carga suspendida ni con manos ocupadas en una actuación insegura. Visitantes, tripulaciones extranjeras y charter usarán todas sus interfaces y comunicaciones en español e inglés. La formación y evaluación incluirán personal nuevo y monitores renovados antes del congelamiento, impartidas por Synaptix, sin depender de un administrador o formador del CLIENTE. Evidencia: tareas/perfiles, idioma, condiciones, incidentes de uso y evaluación individual.

\subsection{CP-20: privacidad, auditoría y firma}
\label{caso:cp-20}
\textbf{Trazabilidad:} C07 RT-11.10, RT-16.09/14. Se prepararán datos protegidos de salud/menores, armadores, tripulantes, contratistas y posición consentida, con autorizaciones y revocaciones. Se exigirá auditoría de cada acceso a salud/menores, posición compartida y deuda, además de modificaciones, conservando quién, qué, cuándo y resultado sin exponer innecesariamente el contenido sensible. Los casos de firma cubrirán contrato/renovación de amarra, autorización del apoderado, recepción/entrega de custodia, acta de izaje y entrega de residuos al gestor. Cada modalidad deberá corresponder a la normativa aplicable y comprobar identidad, integridad, vigencia y conservación. No se aceptará un clic genérico como demostración de firma jurídicamente admisible en todos los casos.

\subsection{CP-21: portales, baja capacidad y notificaciones}
\label{caso:cp-21}
\textbf{Trazabilidad:} C07 RT-16.21/30 y RT-17.01. Se prepararán datos del portal público, armador, apoderado, visitante y contratista y destinos de prueba. Sin autenticación se comprobarán disponibilidad/tarifas del servicio de visita, condiciones de acceso, estado y avisos, en español e inglés. Autenticado se verificarán reserva/pago, contrato/consumo/cuenta/documentos/plan del armador, autorización y salida/regreso escolar, acreditación y residuos del contratista. Se comprobarán avisos de retorno a apoderados y de eventos/cuenta a armadores, más canal de baja capacidad con latencia de horas y mensajes desordenados. Se exigirá correspondencia con la autoridad, aislamiento entre cuentas y un único efecto por solicitud, conservando recibos y estados de entrega.

\subsection{CP-22: servicio gestionado y restricciones de intervención}
\label{caso:cp-22}
\textbf{Trazabilidad:} C07 RT-10.05, RT-15.02, RT-21.06/16 y RT-22.04. Proyecto preparará el calendario real de campañas y 14 eventos, congelamiento del 15 de diciembre al 15 de marzo, avisos de marejada y acceso al campo de 40 boyas; verificará logística marítima segura considerando 6 millas náuticas y 22 km por camino, sin suponer acceso permanente o personal residente. Se comprobarán suspensión inmediata, reprogramación y capacidad de contingencia sin desplazar hitos. El expediente de capacidad incluirá experiencia acreditada en normativa marítima/residuos y servicio gestionado sin TI interna. Se verificarán turnos efectivos de mesa 06:00--24:00 en temporada y 08:00--20:00 el resto, atención en ambos idiomas y emergencia/planes vencidos 24x7x365, con escalamiento y respuesta crítica preservados. Los contratos/cuadrantes y comprobaciones del servicio acreditarán cobertura; un calendario declarado no demuestra personal contratado ni respuesta obtenida.


La Tabla~\ref{tab:t13-complementos} permite localizar las fichas que comprobarán los 27 complementos específicos del capítulo 15 del caso. Su correspondencia registra planificación de verificaciones; no demuestra pruebas ejecutadas ni reemplaza las cláusulas generales del requisito transversal. La recepción exigirá el resultado de cada condición aplicable, con datos, versión y evidencia completos. \parencite[cap.~15, pp.~34--37]{sd9-c07}
\begin{longtable}{L{.17\textwidth}L{.31\textwidth}L{.42\textwidth}}
\caption{Correspondencia de los 27 complementos del caso con las fichas de T-13}\label{tab:t13-complementos}\\\hline
\EncabezadoTabla{ID de T-12} & \EncabezadoTabla{Fichas propuestas} & \EncabezadoTabla{Control específico}\\\hline\endfirsthead
\hline\EncabezadoTabla{ID de T-12} & \EncabezadoTabla{Fichas propuestas} & \EncabezadoTabla{Control específico}\\\hline\endhead
RT-02.12 & CP-13 (\ref{caso:cp-13}) & Parametrización y crecimiento\\\hline
RT-03.10 & CP-01 (\ref{caso:cp-01}), CP-02 (\ref{caso:cp-02}) & Funciones locales y horizontes completos\\\hline
RT-03.13 & CP-03 (\ref{caso:cp-03}) & Reconexión y conciliación\\\hline
RT-03.24 & CP-14 (\ref{caso:cp-14}) & Zonas, marea, segregación y respaldo\\\hline
RT-05.10 & CP-15 (\ref{caso:cp-15}) & Retención por categoría y eliminación\\\hline
RT-05.15 & CP-08 (\ref{caso:cp-08}) & Universo histórico íntegro\\\hline
RT-05.23 & CP-16 (\ref{caso:cp-16}) & Formatos/canales y disponibilidad\\\hline
RT-05.29 & CP-04 (\ref{caso:cp-04}), CP-17 (\ref{caso:cp-17}) & Autoridad y actualización por dominio\\\hline
RT-06.01 & CP-18 (\ref{caso:cp-18}) & Autoadministración y gabinetes\\\hline
RT-09.01 & CP-05 (\ref{caso:cp-05}), CP-06 (\ref{caso:cp-06}) & Tiempos, interacciones y avisos\\\hline
RT-09.02 & CP-17 (\ref{caso:cp-17}) & Volumetría y concentración de peak\\\hline
RT-10.05 & CP-22 (\ref{caso:cp-22}) & Ventanas, eventos y suspensión\\\hline
RT-11.10 & CP-10 (\ref{caso:cp-10}), CP-20 (\ref{caso:cp-20}) & Protección de datos sensibles\\\hline
RT-12.11 & CP-19 (\ref{caso:cp-19}) & Condiciones internas y formación\\\hline
RT-12.12 & CP-19 (\ref{caso:cp-19}), CP-21 (\ref{caso:cp-21}) & Perfiles externos y tareas\\\hline
RT-13.08 & CP-11 (\ref{caso:cp-11}), CP-19 (\ref{caso:cp-19}) & Uso seguro en condiciones de terreno\\\hline
RT-13.12 & CP-19 (\ref{caso:cp-19}), CP-22 (\ref{caso:cp-22}) & Interfaces/comunicaciones y soporte bilingüe\\\hline
RT-15.02 & CP-22 (\ref{caso:cp-22}) & Acreditación sectorial y servicio gestionado\\\hline
RT-16.09 & CP-20 (\ref{caso:cp-20}) & Auditoría de accesos y modificaciones\\\hline
RT-16.14 & CP-20 (\ref{caso:cp-20}) & Modalidad de firma por acto\\\hline
RT-16.21 & CP-06 (\ref{caso:cp-06}), CP-21 (\ref{caso:cp-21}) & Tres canales, acuse, escalamiento y baja capacidad\\\hline
RT-16.30 & CP-21 (\ref{caso:cp-21}) & Portal público y tareas autenticadas\\\hline
RT-17.01 & CP-02 (\ref{caso:cp-02}), CP-19 (\ref{caso:cp-19}), CP-21 (\ref{caso:cp-21}) & Cinco perfiles móviles\\\hline
RT-17.06 & CP-18 (\ref{caso:cp-18}) & Periféricos e integración real\\\hline
RT-21.06 & CP-22 (\ref{caso:cp-22}) & Turnos, idiomas y emergencia permanente\\\hline
RT-21.16 & CP-22 (\ref{caso:cp-22}) & Logística y acceso al campo de boyas\\\hline
RT-22.04 & CP-19 (\ref{caso:cp-19}), CP-22 (\ref{caso:cp-22}) & Renovación estacional y formación antes de peak\\\hline
\end{longtable}
\FuenteTabla{Bases Técnicas del Caso07, cap.~15; IDs conservados en T-12; fichas propuestas por Synaptix.}

Los requisitos que combinan varias condiciones requieren varias comprobaciones. Por ejemplo, RT-16.21 incluye emergencia multicanal y también los avisos operacionales y el canal de baja capacidad; aprobar sólo la emergencia no cierra el requisito. Calidad conservará el desglose por cláusula y sus resultados antes de emitir una decisión global de conformidad.

\section{Defectos, expediente y decisión}
Los defectos conservarán ID, requisito/cláusula, versión, reproducción, severidad, responsable y evidencia. Un crítico/alto o incumplimiento contractual bloqueará promoción; cierre requiere corrección, revisión independiente y repetición del caso y su regresión. El expediente incluirá protocolos, manifiestos, resultados brutos, informes, limitaciones, defectos y decisiones. Se entregará a la Contraparte con T-17, sin convertir cobertura de diseño o existencia de un informe en aceptación. \parencites[sec.~9.2.5--9.2.6]{sd9-sd9}[art.~18.2--18.4, p.~13]{sd9-ba}
